A software quality team that attaches 90\% prediction sets to a defect classifier will usually check the promise on a random test split. In tabular data with repeated rows, this check is misleading. Many test rows have an exact copy in the training data, and the check averages over these easy repeats and the genuinely new cases that matter after deployment. We propose a check--fix--report routine. The check predicts, from the data alone and before any model is trained, how far coverage of new cases can fall. The fix calibrates seen and new cases separately. The report gives new-case coverage next to the usual figure. We tested the routine on \NCsNFlagTotal{} application data sets (software defects, network intrusion, phishing), \NNTotal{} benchmark data sets, controlled copy experiments, and \NScpN{} fresh data sets collected under analysis plans fixed in advance. We used random forests, LightGBM and linear models. On NASA mc1, random-forest sets reported \NCsMcOneReported{}\% coverage but covered \NCsMcOneNew{}\% of new modules; with the fix, \NCsMcOneFixed{}\%. On the fresh data, the check flagged all \NScpPos{} data sets whose new cases fell more than 2 percentage points short (95\% CI for sensitivity \NScpSensCi{}). The fix kept new-case coverage within \NWTwoWorstFixAbs{} points of the target on every data set. The check is a warning and a limit, not a forecast: on the first fresh collection it did not rank data sets by the size of their shortfall. Exact matching also misses near-copies, for which calibration by distance is needed. The routine is available as an open-source tool.
